%! Skills Focus: Selecting, Implementing, and Communicating Inference Procedures — Unit 7, Lesson 7.10 — Exit Ticket
\documentclass[10pt]{article}
\usepackage{apstats-article}
\usepackage{apstats-boxes}
\newcommand{\TallMath}[1]{$\displaystyle #1\rule[-1.4em]{0pt}{3.2em}$}

\begin{document}

\pageheader{Unit 7, Lesson 7.10}{Exit Ticket}
\namedateperiod

\vspace{0.15in}

\textit{For each scenario: (a) name the correct inference procedure, and (b) state one key
condition that must be verified. No computation is required.}

\vspace{0.10in}
\begin{enumerate}

\item A researcher surveys 150 randomly selected registered voters and records whether each
      supports a ballot measure. She wants to \emph{estimate} the true proportion of all
      registered voters who support it.

      \begin{enumerate}[label=\alph*.]
        \item Procedure: \blank{9cm}
        \item One key condition to verify: \writeline
      \end{enumerate}

\vspace{0.12in}

\item A teacher randomly assigns 32 students to one of two tutoring methods and records each
      student's score on a post-test. She wants to test whether the \emph{mean} post-test score
      differs between the two methods.

      \begin{enumerate}[label=\alph*.]
        \item Procedure: \blank{9cm}
        \item One key condition to verify: \writeline
      \end{enumerate}

\vspace{0.12in}

\item A confidence interval for $\mu_1 - \mu_2$ is computed to be $(1.2,\; 4.8)$.\\
      Does the interval provide evidence that $\mu_1 \ne \mu_2$? \blank{2cm}\\
      Explain: \writeline

\end{enumerate}

\end{document}
